% chapters/14/section-03.tex
\section{Canvas: Skor dan Nyawa}
\label{sec:canvas-ui}

\subsection{Mengapa Antarmuka Memerlukan Wadah Tersendiri}

Elemen antarmuka pengguna---teks skor, ikon nyawa, tombol---berbeda sifatnya
dari benda permainan. Ia tidak berada pada ruang permainan, tidak terpengaruh
kamera, dan tidak bergerak ketika pemain berpindah. Karena itu, ia ditempatkan
pada wadah tersendiri yang disebut \istilah{Canvas}.

Susunan antarmuka pada bab ini diperlihatkan pada panel kanan
Gambar~\ref{fig:animator-canvas}. Seluruh elemen berada di
bawah satu \istilah{Canvas}, dan itulah yang membuat susunannya tetap sama
pada setiap ukuran layar \cite{unity-leveldesign}.

\subsection{Tiga Elemen yang Dipakai}

\begin{description}
  \item[\istilah{Text}] Menampilkan tulisan, misalnya skor dan jumlah nyawa.
        Ukuran dan jenis hurufnya diatur pada \istilah{Inspector}.
  \item[\istilah{Image}] Menampilkan gambar, misalnya ikon nyawa atau bingkai
        antarmuka.
  \item[\istilah{Panel}] Wadah untuk mengelompokkan elemen, sehingga
        susunannya dapat dipindahkan bersama-sama.
\end{description}

\subsection{Menghubungkan Antarmuka dengan Keadaan Permainan}

Antarmuka tidak menyimpan keadaan permainan; ia hanya \textbf{menampilkannya}.
Keadaan yang sebenarnya tetap disimpan pada kode Anda, dan antarmuka
diperbarui ketika keadaannya berubah.

\begin{kodekeluaran}
public void TambahSkor(int nilai)
{
    score += nilai;
    scoreText.text = $"Skor: {score}";
}
\end{kodekeluaran}

Cara itu dipilih dengan sengaja. Apabila skor disimpan pada tulisan di layar
saja, maka keadaan permainannya menjadi bergantung pada tampilannya, dan
penyimpanan keadaan menjadi sulit. Aturan yang berlaku sama dengan yang telah
Anda pelajari pada Bagian~3.4: keadaan permainan disimpan pada satu tempat
yang jelas, sedangkan tampilannya hanya mengikuti.

\begin{tipsbox}
Perhatikan pemakaian \code{Mathf.Max(0, lives - nilai)} pada skrip Daftar~\ref{lst:14}: nyawa tidak boleh bernilai negatif, sebab angka negatif pada
antarmuka terlihat janggal. Kebiasaan membatasi nilai pada rentang yang masuk
akal mencegah gejala yang serupa pada banyak hal lain, misalnya skor, waktu,
dan isi persediaan.
\end{tipsbox}

\begin{catatabox}
Rujukan ke elemen antarmuka dihubungkan melalui \istilah{Inspector}, bukan
melalui kode. Apabila rujukan itu belum dihubungkan, program tetap berjalan
tetapi antarmukanya tidak diperbarui, dan kesalahannya muncul sebagai pesan
peringatan pada panel \istilah{Console}. Periksa panel itu setiap kali
antarmuka Anda tidak berubah seperti yang diharapkan.
\end{catatabox}
